\documentclass[11pt,a4paper]{article}

\usepackage[utf8]{inputenc}
\usepackage[T1]{fontenc}
\usepackage[margin=2.5cm]{geometry}
\usepackage{amsmath,amssymb,amsthm}
\usepackage{graphicx}
\usepackage{hyperref}
\usepackage{booktabs}
\usepackage{enumitem}
\usepackage{xcolor}

\title{A Gentle Introduction to Graph Neural Networks}
\author{Jane Researcher \\ \texttt{jane@university.edu}}
\date{\today}

\begin{document}

\maketitle

\begin{abstract}
Graph Neural Networks (GNNs) have emerged as a powerful framework for learning on graph-structured data.
This paper provides an accessible introduction to the core concepts behind GNNs, including message passing,
graph convolutions, and common architectures. We derive the mathematical foundations and discuss
practical considerations for implementation.
\end{abstract}

\section{Introduction}
Graphs are ubiquitous in modern machine learning applications, from social network analysis to molecular
property prediction. Unlike images or sequences, graphs have an irregular structure that does not lend
itself to standard convolutional or recurrent architectures.

A graph $\mathcal{G} = (\mathcal{V}, \mathcal{E})$ consists of a set of nodes $\mathcal{V} = \{v_1, \ldots, v_N\}$
and edges $\mathcal{E} \subseteq \mathcal{V} \times \mathcal{V}$. Each node $v_i$ may have an associated
feature vector $\mathbf{x}_i \in \mathbb{R}^d$.

\section{Message Passing Framework}
The message passing framework \cite{gilmer2017neural} unifies many GNN architectures. At each layer $k$,
every node $v$ computes:

\begin{equation}
\mathbf{h}_v^{(k+1)} = \phi\left(
    \mathbf{h}_v^{(k)},\;
    \bigoplus_{u \in \mathcal{N}(v)} \psi\left(\mathbf{h}_v^{(k)}, \mathbf{h}_u^{(k)}, \mathbf{e}_{vu}\right)
\right)
\label{eq:message_passing}
\end{equation}

where $\mathcal{N}(v)$ denotes the neighbors of $v$, $\psi$ is the \textit{message function},
$\bigoplus$ is a permutation-invariant aggregation operator (sum, mean, or max), and $\phi$ is the
\textit{update function}.

\subsection{Graph Convolutional Networks}
The Graph Convolutional Network (GCN) \cite{kipf2017semi} uses a simple weighted sum aggregation:

\begin{equation}
\mathbf{H}^{(k+1)} = \sigma\left(
    \tilde{\mathbf{D}}^{-\frac{1}{2}} \tilde{\mathbf{A}} \tilde{\mathbf{D}}^{-\frac{1}{2}}
    \mathbf{H}^{(k)} \mathbf{W}^{(k)}
\right)
\label{eq:gcn}
\end{equation}

where $\tilde{\mathbf{A}} = \mathbf{A} + \mathbf{I}$ is the adjacency matrix with self-loops,
$\tilde{\mathbf{D}}_{ii} = \sum_j \tilde{\mathbf{A}}_{ij}$ is the degree matrix,
$\mathbf{W}^{(k)}$ is a learnable weight matrix, and $\sigma$ is a nonlinearity.

\subsection{Graph Attention Networks}
Graph Attention Networks (GATs) \cite{velickovic2018graph} introduce learnable attention coefficients:

\begin{equation}
\alpha_{vu} = \frac{
    \exp\left(\text{LeakyReLU}\left(\mathbf{a}^\top [\mathbf{W}\mathbf{h}_v \,\|\, \mathbf{W}\mathbf{h}_u]\right)\right)
}{
    \sum_{w \in \mathcal{N}(v)} \exp\left(\text{LeakyReLU}\left(\mathbf{a}^\top [\mathbf{W}\mathbf{h}_v \,\|\, \mathbf{W}\mathbf{h}_w]\right)\right)
}
\label{eq:gat_attention}
\end{equation}

The updated node representation is then:

\begin{equation}
\mathbf{h}'_v = \sigma\left(\sum_{u \in \mathcal{N}(v)} \alpha_{vu} \mathbf{W}\mathbf{h}_u\right)
\label{eq:gat_update}
\end{equation}

Multi-head attention extends this by computing $K$ independent attention mechanisms and concatenating
(or averaging) their outputs.

\section{Expressiveness and Limitations}
The expressive power of GNNs is bounded by the 1-dimensional Weisfeiler-Lehman (1-WL) graph isomorphism
test \cite{xu2019powerful}. Two non-isomorphic graphs that 1-WL cannot distinguish will also be
indistinguishable to any standard GNN.

To overcome this limitation, researchers have proposed:
\begin{itemize}[leftmargin=*]
    \item \textbf{Positional encodings:} Inject node position information via Laplacian eigenvectors or random walk features.
    \item \textbf{Higher-order GNNs:} Operate on $k$-tuples of nodes rather than individual nodes.
    \item \textbf{Subgraph GNNs:} Process local subgraphs to break symmetry.
\end{itemize}

\begin{table}[htbp]
    \centering
    \caption{Comparison of GNN architectures}
    \label{tab:gnn_comparison}
    \begin{tabular}{@{}lccc@{}}
        \toprule
        Architecture & Aggregation & Attention & Complexity \\
        \midrule
        GCN          & Mean        & No        & $\mathcal{O}(|\mathcal{E}|)$ \\
        GraphSAGE    & Mean/Max    & No        & $\mathcal{O}(|\mathcal{E}|)$ \\
        GAT          & Weighted Sum& Yes       & $\mathcal{O}(|\mathcal{E}|)$ \\
        GIN          & Sum         & No        & $\mathcal{O}(|\mathcal{E}|)$ \\
        MPNN         & Any         & Optional  & Varies \\
        \bottomrule
    \end{tabular}
\end{table}

\section{Empirical Evaluation}
We evaluate the architectures on three standard benchmarks: Cora, CiteSeer, and PubMed for node
classification, and PROTEINS and NCI1 for graph classification. Table~\ref{tab:results} summarizes
our findings.

\begin{table}[htbp]
    \centering
    \caption{Test accuracy (\%) across benchmarks}
    \label{tab:results}
    \begin{tabular}{@{}lccccc@{}}
        \toprule
        Model    & Cora  & CiteSeer & PubMed & PROTEINS & NCI1 \\
        \midrule
        GCN      & 81.5  & 70.3     & 79.0   & 76.0     & 80.2 \\
        GAT      & 83.0  & 72.5     & 79.0   & 74.7     & 78.5 \\
        GIN      & 82.1  & 71.8     & 78.8   & 76.2     & 82.7 \\
        GraphSAGE& 81.8  & 71.2     & 78.5   & 75.9     & 77.7 \\
        \bottomrule
    \end{tabular}
\end{table}

\section{Conclusion}
We have provided a mathematical introduction to graph neural networks, covering the message passing
paradigm, major architectural variants, and theoretical limitations. GNNs represent a powerful and
flexible framework that continues to drive advances across diverse scientific domains.

\bibliographystyle{plain}
\begin{thebibliography}{9}

\bibitem{gilmer2017neural}
J.~Gilmer, S.~Schoenholz, P.~Riley, O.~Vinyals, and G.~Dahl.
\newblock Neural message passing for quantum chemistry.
\newblock In \emph{ICML}, 2017.

\bibitem{kipf2017semi}
T.~Kipf and M.~Welling.
\newblock Semi-supervised classification with graph convolutional networks.
\newblock In \emph{ICLR}, 2017.

\bibitem{velickovic2018graph}
P.~Veli\v{c}kovi\'{c}, G.~Cucurull, A.~Casanova, A.~Romero, P.~Li\`{o}, and Y.~Bengio.
\newblock Graph attention networks.
\newblock In \emph{ICLR}, 2018.

\bibitem{xu2019powerful}
K.~Xu, W.~Hu, J.~Leskovec, and S.~Jegelka.
\newblock How powerful are graph neural networks?
\newblock In \emph{ICLR}, 2019.

\end{thebibliography}

\end{document}
